\documentclass[12pt]{article}

\usepackage{fontawesome}
\usepackage{hyperref}
\usepackage{xurl}
\usepackage[tagged, highstructure]{accessibility}
\usepackage{bookmark}
\usepackage{enumitem}

\hypersetup{
    colorlinks=false,
    pdfborder={0 0 0},
    pdftitle={Big Data Lab},
    pdfauthor={Adrianna Holden-Gouveia},
    pdfsubject={Introduction to Data - Lab Assignment},
    pdflang={en-US},
    pdfkeywords={big data, machine learning, Kaggle, Titanic dataset, predictive analytics}
}

% Configure PDF bookmarks for navigation
\bookmarksetup{
    numbered,
    open,
}

% Configure list spacing for better accessibility
\setlist{nosep}



\title{Big Data}
\author{
        Adrianna Holden-Gouveia \\
        Website: \href{https://aholdengouveia.name}{aholdengouveia.name}\\ 
        \date{\vspace{-5ex}}
        %Email: \href{mailto:admin@aholdengouveia.name}{admin@aholdengouveia.name} \\
%        \faLinkedin{: aholdengouveia} \\
        \faGithub {: aholdengouveia} \\
        %\faTwitter {: aholdengouveia} \\
        }

%S\date{\today}


\begin{document}    

\maketitle

%\begin{abstract}
%This is a template for Linux Administration Lab work
%\end{abstract}
%\tableofcontents


\section*{Accessibility Notice}
This document is also available in HTML format at:

Link: \href{https://aholdengouveia.name/IntroData/labs/bigdata.html}{\textbf{\nolinkurl{https://aholdengouveia.name/IntroData/labs/bigdata.html}}}

The HTML version provides enhanced accessibility features including keyboard navigation, screen reader support, responsive design, dark mode support, and high contrast options.

\section*{Objectives:}
\begin{itemize}
    \item Students will gain a basic understanding of how to work through a problem using a larger dataset and some basic Machine Learning. 
\end{itemize}
\section*{Complete the following}

References, a video, a PowerPoint and some notes are available at my website
Link: \href{https://www.aholdengouveia.name/IntroData/bigdata.html}{\textbf{\nolinkurl{https://www.aholdengouveia.name/IntroData/bigdata.html}}}

\subsection*{Steps}
\begin{enumerate}
    \item Go to Kaggle Link: \href{https://www.kaggle.com/}{\textbf{\nolinkurl{https://www.kaggle.com/}}} and create a free account.
    \item Join the Titanic competition. The tutorial below will walk you through this.
    \item Complete the Titanic tutorial. Link: \href{https://www.kaggle.com/code/alexisbcook/titanic-tutorial}{\textbf{\nolinkurl{https://www.kaggle.com/code/alexisbcook/titanic-tutorial}}}
    \item Submit your predictions to the competition.
    \item Take a screenshot of your score on the leaderboard. Your screenshot must show your Kaggle username, and include a note with your name and the term, the same as the screenshots in your other labs.
\end{enumerate}

\subsection*{Please include Answers to the following questions}
    \begin{enumerate}
        \item What do you think of the tutorial?
        \item How difficult did you find this activity? Why?
        \item Will you be trying another challenge on Kaggle? Why or why not? If yes, which one?
        \item Improving the code
        \begin{enumerate}
            \item Were you able to improve on the given code?
            \item What did you try?
            \item What was your score before and after?
        \end{enumerate}
        \item What was your accuracy? You can find this on your Submissions page, where it might show up as \textbf{Score} or \textbf{Public Score}.
        \item Do you have any ideas on how you could improve that?
    \end{enumerate}

\section*{Deliverables}
\begin{enumerate}
    \item A text document with the answers to the questions.
    \item A screenshot of your score on the leaderboard, showing your Kaggle username and a note with your name and the term.
\end{enumerate}
\end{document}